% Compile with XeLaTeX or LuaLaTeX
\documentclass[14pt,a4paper]{extreport}

% -------------------------
% Language and font setup
% -------------------------
\usepackage{fontspec}
\usepackage{polyglossia}
\setmainlanguage{russian}
\setotherlanguage{english}

\setmainfont{Times New Roman}

% -------------------------
% Page layout
% -------------------------
\usepackage[
    a4paper,
    left=25mm,   % left margin >= 20mm
    right=15mm,
    top=20mm,
    bottom=20mm
]{geometry}

% -------------------------
% Paragraph formatting
% -------------------------
\usepackage{setspace}
\setstretch{1.5}                 % allowed range: 1.3--1.5
\setlength{\parindent}{1cm}      % red line / paragraph indentation
\setlength{\parskip}{0pt}

\usepackage{ragged2e}
\justifying

% -------------------------
% Math, figures, tables
% -------------------------
\usepackage{amsmath, amssymb}
\usepackage{graphicx}
\usepackage{float}
\usepackage{booktabs}
\usepackage{array}
\usepackage{longtable}
\usepackage{caption}
\usepackage{chngcntr}

% Number figures, tables, and equations inside chapters: 2.3, 2.4, etc.
\numberwithin{figure}{chapter}
\numberwithin{table}{chapter}
\numberwithin{equation}{chapter}

% Figure captions: below, centered, "Рис. 2.3."
\captionsetup[figure]{
    name=Рис.,
    labelsep=period,
    justification=centering,
    singlelinecheck=false
}

% Table captions: above, right-aligned, "Таблица 2.3."
\captionsetup[table]{
    name=Таблица,
    labelsep=period,
    justification=raggedleft,
    singlelinecheck=false,
    position=top
}

% -------------------------
% Section heading formatting
% -------------------------
\usepackage{titlesec}

\titleformat{\chapter}
    {\normalfont\bfseries\centering}
    {\thechapter}
    {1em}
    {}

\titleformat{\section}
    {\normalfont\bfseries}
    {\thesection}
    {1em}
    {}

\titleformat{\subsection}
    {\normalfont\bfseries}
    {\thesubsection}
    {1em}
    {}

% -------------------------
% Table of contents formatting
% -------------------------
\usepackage{tocloft}
\renewcommand{\contentsname}{Содержание}

% -------------------------
% Bibliography: GOST-style references
% -------------------------
% Requires biblatex-gost package.
% Compile sequence:
% xelatex main.tex
% biber main
% xelatex main.tex
% xelatex main.tex
\usepackage[
    backend=biber,
    style=gost-numeric,
    sorting=none
]{biblatex}

\addbibresource{references.bib}

% -------------------------
% Hyperlinks and cross-references
% -------------------------
\usepackage[hidelinks]{hyperref}

% -------------------------
% Appendix formatting
% -------------------------
\usepackage{appendix}

\renewcommand{\appendixname}{Приложение}
\renewcommand{\appendixtocname}{Приложения}
\renewcommand{\appendixpagename}{Приложения}

% -------------------------
% Helper commands
% -------------------------

% Unnumbered annotation page
\newcommand{\annotationpage}[2]{
    \clearpage
    \thispagestyle{empty}
    \begin{center}
        \textbf{#1}
    \end{center}

    #2
}

% Standard chapter conclusion section
\newcommand{\chapterconclusions}{
    \section{Выводы и результаты по главе}
}

% Formula legend helper
\newenvironment{formulalegend}{
    \par\noindent\textit{где:}
    \begin{itemize}
}{
    \end{itemize}
}

% -------------------------
% Document begins
% -------------------------
\begin{document}

% =========================================================
% Title page: not numbered
% =========================================================
\pagenumbering{gobble}

\begin{titlepage}
    \begin{center}
        Федеральное государственное автономное образовательное учреждение\\
        высшего образования\\
        «Национальный исследовательский университет\\
        “Высшая школа экономики”»\\[1.5cm]

        Факультет Компьютерных Наук\\
        Образовательная программа «Современные Компьютерные Науки»\\[2.5cm]

        \textbf{ВЫПУСКНАЯ КВАЛИФИКАЦИОННАЯ РАБОТА}\\[0.5cm]

        на тему:\\[0.5cm]

        \textbf{Квантизация и сжатие генеративных диффузионных моделей}\\[2cm]
    \end{center}

    \begin{flushright}
        Выполнил студент группы мСКН241:\\
        Пыж Владислав Александрович\\[0.5cm]

        Научный руководитель:\\
        канд. наук, доцент\\
        Аланов Айбек
    \end{flushright}

    \vfill

    \begin{center}
        Москва, 2026
    \end{center}
\end{titlepage}

% =========================================================
% Russian annotation: not numbered
% Requirement: 1500--2000 characters including spaces
% =========================================================
\annotationpage{Аннотация}{
Аннотация представляет собой краткую характеристику выпускной квалификационной работы.
В ней необходимо описать цель и задачи исследования, объект и предмет работы,
используемые методы, эмпирическую базу при наличии прикладного исследования,
а также основные результаты, полученные в ходе выполнения работы.

Текст данной аннотации следует заменить на собственный текст объемом от 1500
до 2000 знаков с учетом пробелов. Аннотация должна быть написана на русском
языке и размещена на отдельном листе после титульного листа.
}

% =========================================================
% English annotation: not numbered
% Requirement: 1500--2000 characters including spaces
% =========================================================
\annotationpage{\textenglish{Abstract}}{
\begin{english}
The abstract is a concise summary of the graduation thesis. It should describe
the purpose and objectives of the research, the object and subject of the study,
the methodology, the empirical basis if the work includes applied research,
and the main results obtained in the thesis.

Replace this placeholder with your own English abstract. The required length is
from 1500 to 2000 characters including spaces. The English abstract should be
placed on a separate page after the Russian annotation.
\end{english}
}

% =========================================================
% Main part starts here
% =========================================================
\clearpage
\pagenumbering{arabic}

\tableofcontents
\clearpage

% =========================================================
% Introduction
% =========================================================
\chapter*{Введение}
\addcontentsline{toc}{chapter}{Введение}

Во введении необходимо раскрыть тему работы, объект, предмет и методы
исследования, кратко обосновать актуальность и значимость темы, сформулировать
цель и основные задачи, описать основной результат и структуру работы.

Пример ссылки на источник: \cite{example-source}.
Пример перекрестной ссылки: см. рисунок~\ref{fig:example} и таблицу~\ref{tab:example}.

% =========================================================
% Chapter 1
% =========================================================
\chapter{Название первой главы}

Текст первой главы. Заголовки разделов не должны содержать сокращений и
аббревиатур, за исключением общепринятых.

\section{Название раздела}

Пример рисунка приведен ниже.

\begin{figure}[H]
    \centering
    \includegraphics[width=0.65\textwidth]{example-image}
    \caption{Пример визуализации полученных данных}
    \label{fig:example}
\end{figure}

Пример таблицы приведен ниже.

\begin{table}[H]
    \caption{Объемные характеристики программных модулей}
    \label{tab:example}
    \centering
    \begin{tabular}{lcc}
        \toprule
        Модуль & Количество строк & Количество функций \\
        \midrule
        Модуль A & 120 & 8 \\
        Модуль B & 250 & 15 \\
        Модуль C & 180 & 11 \\
        \bottomrule
    \end{tabular}
\end{table}

Пример выделенной формулы:

\begin{equation}
    y = f(x; \theta)
    \label{eq:model}
\end{equation}

\begin{formulalegend}
    \item $x$ --- входные данные;
    \item $y$ --- результат работы модели;
    \item $\theta$ --- параметры модели;
    \item $f$ --- функция преобразования входных данных.
\end{formulalegend}

\chapterconclusions

В данном разделе необходимо кратко изложить результаты главы и обеспечить
логический переход к следующей главе.

% =========================================================
% Chapter 2
% =========================================================
\chapter{Название второй главы}

Текст второй главы.

\section{Пример длинной таблицы}

Если таблица занимает более одной страницы, можно использовать окружение
\texttt{longtable}. Подпись будет повторяться на следующих страницах с пометкой
«продолжение».

\begin{longtable}{p{0.25\textwidth}p{0.6\textwidth}}
    \caption{Пример длинной таблицы}
    \label{tab:long-example}
    \\
    \toprule
    Параметр & Описание \\
    \midrule
    \endfirsthead

    \caption*{\raggedleft Таблица \thetable. Пример длинной таблицы (продолжение)}
    \\
    \toprule
    Параметр & Описание \\
    \midrule
    \endhead

    \bottomrule
    \endfoot

    Параметр 1 & Описание параметра 1. \\
    Параметр 2 & Описание параметра 2. \\
    Параметр 3 & Описание параметра 3. \\
    Параметр 4 & Описание параметра 4. \\
    Параметр 5 & Описание параметра 5. \\
\end{longtable}

\chapterconclusions

Краткие выводы и результаты по второй главе.

% =========================================================
% Conclusion
% =========================================================
\chapter*{Заключение}
\addcontentsline{toc}{chapter}{Заключение}

В заключении необходимо кратко описать полученные результаты, степень достижения
цели и решения поставленных задач, а также перспективы дальнейшей работы.

% =========================================================
% Bibliography
% =========================================================
\printbibliography[title={Библиографический список}]

% =========================================================
% Appendices
% =========================================================
\begin{appendices}

\chapter{Терминологический словарь}

Пример приложения. Приложения должны быть пронумерованы и перечислены в
содержании работы.

\chapter{Список сокращений}

\begin{description}
    \item[ВКР] выпускная квалификационная работа;
    \item[НИУ ВШЭ] Национальный исследовательский университет «Высшая школа экономики».
\end{description}

\chapter{Ключевые фрагменты исходного кода}

Не рекомендуется включать полный исходный код всех разработанных программ.
Можно включать только ключевые фрагменты, необходимые для демонстрации
оригинальных решений.

\end{appendices}

\end{document}